%!TEX TS-program = xelatex
\documentclass[12pt,a4paper]{ctexart}

\usepackage{graphicx}
\usepackage{booktabs}
\usepackage{longtable}
\usepackage{hyperref}
\usepackage[margin=2.5cm]{geometry}
\usepackage{titlesec}
\usepackage{xcolor}
\usepackage{fancyhdr}

\hypersetup{
    colorlinks=true,
    linkcolor=blue,
    urlcolor=blue,
    pdftitle={第一阶段数据收集与整理报告},
    pdfauthor={Qwen3-8B微调项目},
}

\pagestyle{fancy}
\fancyhf{}
\fancyhead[L]{Qwen3-8B 微调项目}
\fancyhead[R]{第一阶段报告}
\fancyfoot[C]{\thepage}

\title{\textbf{大模型微调项目第一阶段报告}\\[4pt]
       \large 数据来源、采集与整理}
\author{计算机·密码与安全特化模型}
\date{\today}

\begin{document}

\maketitle
\thispagestyle{empty}

\begin{abstract}
\noindent 本报告记录了基于 Qwen3-8B 的计算机、密码学与安全领域大模型微调项目的第一阶段工作——数据来源、采集与整理。项目通过多源爬取、PDF教材处理（原生文字提取与OCR）等方式，共采集处理 63 本 PDF、804 篇百科词条、60 篇 RFC 文档、30 篇 OWASP 安全指南及 24 条 NIST 存根，最终产出 19,464 篇结构化原始文本块（29.6 MB），统一以最简 JSONL 格式（\texttt{\{"text": "..."\}}）输出，可直接用于后续问答对构建与指令微调训练。
\end{abstract}

\section{项目概述}

本项目以本地部署的 Qwen3-8B 作为基础模型，聚焦于计算机科学、密码学与网络安全三大领域，通过领域数据收集、文本提取、结构化整理，构建面向教材解读与专业知识问答的微调数据集。

第一阶段的\textbf{核心目标}是：从多种来源收集高质量原始语料，统一清洗、分块、格式化，为后续指令数据构建和 LoRA/QLoRA 微调提供结构化的训练基础。

\section{数据来源总览}

\subsection{来源分类}

数据来源分为三大类：\textbf{公开网络资源}（主动爬取）、\textbf{权威标准文档}（下载与提取）、\textbf{专业教材}（PDF 处理）。

\begin{table}[htbp]
\centering
\caption{数据来源汇总}
\label{tab:sources}
\begin{tabular}{@{}llrrr@{}}
\toprule
\textbf{类别} & \textbf{来源} & \textbf{数量} & \textbf{大小} & \textbf{方式} \\
\midrule
公开网络 & Wikipedia (中英) & 804 篇 & 0.95 MB & API 爬取 \\
         & RFC Editor       & 60 篇  & 5.94 MB & 直接下载 \\
         & OWASP            & 30 篇  & 0.87 MB & GitHub Raw \\
         & NIST CSRC        & 24 条  & 0.01 MB & API 爬取(存根) \\
\midrule
权威标准 & NIST SP 800 系列  & 30 本  & ---    & 直链下载+提取 \\
         & NIST FIPS 系列   & 9 本   & ---    & 直链下载+提取 \\
\midrule
专业教材 & 计算机科学教材   & 7 本   & 3,865 chunks & PDF 原生/OCR \\
         & 密码学教材       & 7 本   & 3,495 chunks & PDF 原生/OCR \\
         & 安全类教材       & 9 本   & 3,775 chunks & PDF 原生/OCR \\
\bottomrule
\end{tabular}
\end{table}

\subsection{公开网络资源详情}

\subsubsection*{Wikipedia}

通过 Wikipedia API 抓取英文和中文两个语言的计算机科学、密码学、安全三大领域词条。每个领域下设若干子分类（如 algorithm, data\_structure, cryptography, computer\_security 等），每个分类最多抓取 50 篇词条。使用 API 的 extract 功能获取纯文本摘要。

\begin{itemize}
    \item 英文词条：覆盖 Computer Science (12 分类)、Cryptography (11 分类)、Security (13 分类)
    \item 中文词条：覆盖 计算机科学、密码学、网络安全等 10 个分类
    \item 平均长度：$\sim$950 字符/篇
\end{itemize}

\subsubsection*{RFC Editor}

下载与密码学和安全相关的 RFC 文档全文。包括手动精选的 60 篇核心 RFC（TLS、PKI、SSH、IPsec、DNSSEC、OAuth、JWT 等）和关键词搜索补充。RFC 文档为完整协议规范，平均长度 $\sim$100,000 字符/篇，是数据集中最厚重的单篇文档。

\subsubsection*{OWASP}

从 OWASP GitHub 仓库获取 Cheat Sheet Series（32 篇）、Application Security Verification Standard (ASVS) 章节、Top 10 风险类别等安全实践指南。内容涵盖认证、授权、输入验证、XSS 防护、SQL 注入防护、密码存储、TLS 配置等。

\subsubsection*{NIST CSRC}

NIST 计算机安全资源中心（CSRC）网站为 React SPA 架构，自动化抓取受阻。当前仅通过 API 获取了 24 条出版物元数据存根（SP 800 系列 14 条、FIPS 系列 10 条）。完整内容通过 PDF 下载方式另行获取（见下节）。

\subsection{权威标准文档——NIST PDF}

为弥补 NIST 网站爬取困难的问题，直接从 \texttt{nvlpubs.nist.gov} 下载了 39 篇 NIST 出版物的 PDF 文件。这些均为英文原生文字 PDF，使用 pymupdf4llm 直接提取 markdown 后分块。

\begin{itemize}
    \item SP 800 系列（30 篇）：涵盖安全管理(53, 37)、密码学(38, 52, 56, 57, 90, 108, 131, 133, 175, 186)、网络安全(77, 161, 171, 190, 204, 207)、身份管理(63)
    \item FIPS 系列（9 篇）：140-3(密码模块)、180-4(SHS)、186-5(DSS)、197(AES)、198-1(HMAC)、199(安全分类)、200(最低安全要求)、201-3(PIV)、202(SHA-3)
\end{itemize}

\subsection{专业教材 PDF}

通过 mapping\_knowledge\_domain 子项目的 pdf\_slicer.py 工具处理教材 PDF。根据 PDF 的文字层特征分为两条处理路径：

\begin{description}
    \item[原生文字 PDF（5 本）] 使用 pymupdf4llm 直接提取结构化 markdown，速度快、质量高。包括：公钥密码学的数学基础（王小云）、算法数论、云计算安全技术与应用、云计算数据安全、网络空间安全战略。
    \item[扫描版 PDF（18 本）] 使用 PyMuPDF 逐页渲染为图像，通过 Tesseract OCR（chi\_sim+eng）识别中文文字，经 CJK 字间空格修复和低内容页面过滤后分块。包括：王道408考研系列(4本)、TCP/IP详解、TAOCP、密码编码学与网络安全(Stallings)、应用密码学(Schneier)、深入浅出密码学(Paar)、计算机安全学(Gollmann) 等。
\end{description}

\section{数据处理流程}

\subsection{整体架构}

数据处理分为\textbf{三个主要管道}，最终汇聚为统一的 raw JSONL 格式：

\begin{enumerate}
    \item \textbf{爬虫管道}（wikipedia/rfc/owasp/nist）：API请求/直接下载 $\to$ HTML/Markdown清洗 $\to$ JSONL存储
    \item \textbf{英文PDF管道}（NIST/原生英文）：pymupdf4llm提取 $\to$ 英文清洗(去连字符断行/页眉页脚) $\to$ 语义分块 $\to$ raw JSONL
    \item \textbf{中文PDF管道}（教材/扫描版）：PyMuPDF渲染 $\to$ Tesseract OCR $\to$ CJK去空格 $\to$ 低内容页过滤 $\to$ 语义分块 $\to$ raw JSONL
\end{enumerate}

\subsection{文本清洗}

\begin{itemize}
    \item Wikipedia：去除引用标记 [N]、编辑标记 [edit]、目录魔法字 \_\_TOC\_\_，截断 References/See also/External links 章节
    \item RFC：去除页眉（Network Working Group等）、页脚（[Page N]）、状态栏
    \item OWASP：去除 YAML frontmatter 和 HTML 注释
    \item 英文PDF：去连字符断行（hyphen- ation $\to$ hyphenation）、合并软换行、去独立页码行
    \item 中文OCR：循环去除CJK字符间空格、标点前空格清理
\end{itemize}

\subsection{语义分块}

分块策略：优先按 \#\# 标题边界分割，其次按段落（空行），再次按句子边界（。！？.!?）。块大小上限 1500 字符，相邻块 150 字符重叠（重叠用于上下文连贯，训练时可通过 deoverlap 去除）。

\subsection{分类整理}

通过基于关键词匹配加来源加权的分类器，将所有文档归入三个领域。密码学与安全合并为同一大类（\texttt{crypto\_and\_security}），因为安全架构本身包含密码实现。

\section{最终数据集}

\subsection{数据统计}

\begin{table}[htbp]
\centering
\caption{最终 raw 数据集统计}
\label{tab:final}
\begin{tabular}{@{}lrrr@{}}
\toprule
\textbf{领域} & \textbf{文档数} & \textbf{大小(MB)} & \textbf{子类数} \\
\midrule
computer\_science   & 6,221 & 8.3  & 8 \\
crypto\_and\_security & 13,026 & 17.5 & 9 \\
comprehensive       & 217   & 3.8  & 2 \\
\midrule
\textbf{总计}       & \textbf{19,464} & \textbf{29.6} & \textbf{19} \\
\bottomrule
\end{tabular}
\end{table}

\begin{table}[htbp]
\centering
\caption{各子类分布}
\label{tab:subcats}
\small
\begin{tabular}{@{}llr@{}}
\toprule
\textbf{领域} & \textbf{子类} & \textbf{文档数} \\
\midrule
\multirow{8}{*}{computer\_science}
    & algorithms\_and\_data\_structures  & 1,187 \\
    & computer\_networks                  & 648 \\
    & operating\_systems                  & 121 \\
    & theory\_and\_discrete\_math          & 69 \\
    & databases                           & 39 \\
    & programming\_languages\_and\_se       & 38 \\
    & computer\_architecture              & 20 \\
    & artificial\_intelligence            & 13 \\
    & (OCR教材chunks)                     & $\sim$4,086 \\
\midrule
\multirow{9}{*}{crypto\_and\_security}
    & key\_management\_and\_pki            & 2,992 \\
    & cryptographic\_algorithms           & 1,289 \\
    & cryptographic\_foundations          & 1,229 \\
    & network\_security                   & 525 \\
    & authentication\_and\_access\_control & 503 \\
    & security\_engineering               & 394 \\
    & system\_security                    & 189 \\
    & cryptographic\_protocols            & 93 \\
    & web\_security                       & 53 \\
    & (OCR教材chunks)                     & $\sim$5,759 \\
\midrule
\multirow{2}{*}{comprehensive}
    & general                             & 163 \\
    & cross\_domain                        & 55 \\
\bottomrule
\end{tabular}
\end{table}

\subsection{输出格式}

最终输出的训练数据采用\textbf{最简 JSONL 格式}，每行一个 JSON 对象，仅包含 \texttt{text} 字段：

\begin{verbatim}
{"text": "云计算环境是计算与存储并存的网络环境..."}
{"text": "The Advanced Encryption Standard (AES) specifies..."}
{"text": "公钥密码学的安全性基于数学困难问题..."}
\end{verbatim}

该格式相比完整元数据格式节省约 25\% 的存储空间和 token 开销，每行仅 12 字符的 JSON 包装。可直接通过 HuggingFace \texttt{datasets} 库加载：

\begin{verbatim}
from datasets import load_dataset
ds = load_dataset("json", data_files="data/raw_for_training/
                  crypto_and_security/data.jsonl")
\end{verbatim}

\subsection{数据目录结构}

\begin{verbatim}
data/raw_for_training/
├── computer_science/
│   ├── algorithms_and_data_structures.jsonl
│   ├── computer_networks.jsonl
│   ├── operating_systems.jsonl
│   ├── ...
│   └── data.jsonl              ← OCR教材合并
├── crypto_and_security/
│   ├── cryptographic_algorithms.jsonl
│   ├── cryptographic_foundations.jsonl
│   ├── key_management_and_pki.jsonl
│   ├── network_security.jsonl
│   ├── ...
│   └── data.jsonl              ← OCR教材合并
└── comprehensive/
    ├── general.jsonl
    └── cross_domain.jsonl
\end{verbatim}

\section{处理的 PDF 教材清单}

\subsection{计算机科学（7 本）}

\begin{itemize}
    \item 2025王道408-数据结构（咸鱼学长）
    \item 2025王道考研408-计算机组成原理（咸鱼学长）
    \item 2025王道考研408-计算机网络（咸鱼学长）
    \item 2025年操作系统考研复习指导
    \item TCP/IP详解 卷1：协议（原书第2版，Kevin R. Fall）
    \item 计算机概论：基础科学、软件与信息安全导向（北极星）
    \item 计算机程序设计艺术 第1卷 基本算法（Donald E. Knuth）
\end{itemize}

\subsection{密码学与安全（17 本）}

\begin{itemize}
    \item 密码编码学与网络安全 原理与实践 第七版（William Stallings）
    \item 应用密码学：协议、算法与C源程序 原书第2版（Bruce Schneier）
    \item 深入浅出密码学——常见加密技术原理与应用（Christof Paar \& Jan Pelzl）
    \item 公钥密码学的数学基础 第二版（王小云 王明强 孟宪萌）
    \item 算法数论：格、数域、曲线和密码学（J. P. Buhler 等编）
    \item 计算机安全学（Dieter Gollmann）
    \item 网络安全基础：应用与标准 第5版（William Stallings）
    \item 软件加密与解密（Christian Collberg）
    \item 逆向工程实战（Bruce Dang 等）
    \item 逆向工程核心原理（李承远）
    \item 人工智能安全（方滨兴 编）
    \item 云计算安全技术（徐涛 等）
    \item 云计算安全技术与应用（中国电信网络安全实验室）
    \item 云计算数据安全（黄勤龙）
    \item 云存储安全：大数据分析与计算的基石（陈兰香）
    \item 网络空间安全战略（郭宏生 编著）
\end{itemize}

\subsection{NIST 出版物（39 篇）}

包括 SP 800 系列 30 篇和 FIPS 系列 9 篇，涵盖安全管理、密码学标准、网络安全、身份与访问控制等（完整列表见附录或项目目录）。

\section{工具链与脚本清单}

\begin{table}[htbp]
\centering
\caption{项目脚本清单}
\label{tab:scripts}
\small
\begin{tabular}{@{}ll@{}}
\toprule
\textbf{脚本} & \textbf{功能} \\
\midrule
\texttt{main.py}                 & 爬虫主入口，支持按来源运行 \\
\texttt{classify.py}             & 关键词+来源加权分类整理 \\
\texttt{convert\_to\_raw.py}     & curated $\to$ raw 格式转换 \\
\texttt{chunk\_to\_jsonl.py}     & 统一分块→JSONL模块 \\
\texttt{download\_nist\_pdfs.py} & NIST PDF 下载器 \\
\texttt{organize\_pdfs.py}       & 已处理PDF移入\_processed \\
\texttt{tesseract\_ocr.py}       & Tesseract OCR管道(CJK修复+页面过滤) \\
\texttt{process\_english\_pdfs.py} & 英文PDF→raw JSONL \\
\texttt{process\_scanned\_pdfs.py} & 扫描版PDF OCR→raw JSONL \\
\texttt{process\_native\_pdfs.py}  & 原生PDF→raw JSONL \\
\texttt{ocr\_raw\_robust.py}       & 鲁棒OCR→raw(单页容错) \\
\bottomrule
\end{tabular}
\end{table}

外部依赖：pymupdf4llm（PDF结构提取）、Tesseract OCR + chi\_sim（中文识别）、mwclient/wikipedia-api（百科爬取）、beautifulsoup4/html2text（HTML清洗）。

\section{总结与下一步}

\subsection{第一阶段成果}

\begin{itemize}
    \item 建立了完整的多源数据采集管道（爬虫 + PDF提取 + OCR）
    \item 收集处理 63 本专业 PDF 教材和标准文档
    \item 产出 \textbf{19,464 篇} 结构化原始文本，\textbf{29.6 MB}，零格式错误
    \item 全部数据以统一 \texttt{\{"text": "..."\}} JSONL 格式存放于 \texttt{data/raw\_for\_training/}
\end{itemize}

\subsection{第二阶段展望}

第二阶段将基于本阶段产出的 raw 数据，构建面向学术学习场景的高质量指令数据集：

\begin{enumerate}
    \item 从 raw 文本块生成问答对（概念解释、章节总结、方法对比、结论归纳等）
    \item 使用强模型（API）辅助生成 + 人工筛选标注
    \item 构建训练集/验证集/测试集，题型覆盖术语解释、段落总结、论文摘要改写等
    \item 采用 QLoRA 对 Qwen3-8B 进行参数高效微调
    \item 从准确性、完整性、专业性、忠实性四个维度评估微调效果
\end{enumerate}

\end{document}
